% ===========================================================================
% WaterPaper 内置兜底骨架
%
% 使用场景：用户提供的 LaTeX 模板残缺（无 \begin{document} / 无正文注入点 /
% 原样编译失败无法定位），或用户只给了文字格式要求。此时用本骨架 + 从模板
% 抽出的排版参数覆盖）仍能产出 PDF。
%
% 排版对齐 references/default_format.md（docx 默认格式的 LaTeX 等价物）：
%   论文标题  黑体二号加粗居中        \zihao{2}
%   一级标题  黑体小三加粗            \zihao{-3}
%   二级标题  黑体四号加粗            \zihao{4}
%   三级标题  黑体小四加粗            \zihao{-4}
%   正文      宋体小四，1.5 倍行距，首行缩进 2 字符
%   摘要标题  黑体小四加粗居中
%   关键词    标签黑体小四加粗顶格，内容宋体小四，分号分隔
%   图题      宋体五号居中，置于图下方
%   表题      宋体五号加粗居中，置于表上方
%   参考文献  宋体五号（GB/T 7714）
%
% 占位符由 tools/build_paper_pdf.py 替换。注意：本注释区**不得**写出占位符
% 字面量 —— 否则朴素字符串替换会先命中注释，把正文注入到 \documentclass
% 之前，报错还是让人摸不着头脑的 "Missing \begin{document}"。
% build 脚本会断言每个占位符恰好出现一次。
%
% 需替换的占位符（括号内为用途）：
%   WP_TITLE             论文标题
%   WP_AUTHOR            作者（无则留空）
%   WP_STYLE_OVERRIDES   从用户模板抽取的排版参数覆盖（可为空）
%   WP_ABSTRACT_BEGIN / WP_ABSTRACT_END   之间的块在无摘要时整段删除
%   WP_ABSTRACT          摘要正文
%   WP_KEYWORDS          关键词
%   WP_CONTENT           正文（body.tex 内容）
%   WP_BIBLIOGRAPHY      参考文献机制（bibtex 模式下填 \bibliography，
%                        否则留空，因为 thebibliography 已内嵌在正文里）
% ===========================================================================
\documentclass[12pt,a4paper]{ctexart}

\usepackage[top=2.54cm,bottom=2.54cm,left=3.18cm,right=3.18cm]{geometry}
\usepackage{graphicx}
\usepackage{booktabs}
\usepackage{tabularx}
\usepackage{amsmath}
\usepackage{caption}
\usepackage[hidelinks]{hyperref}

% 关键词命令（ctexart 未内置，学校模板通常也自己定义）
\newcommand{\keywords}[1]{%
  \par\vspace{0.8em}\noindent{\heiti\zihao{-4}关键词：}{\zihao{-4}#1}%
}

% 章节层级格式：对齐 default_format.md
\ctexset{
  section = {
    format = \zihao{-3}\heiti,
    name = {,、},
    number = \chinese{section},
    beforeskip = 12pt,
    afterskip = 6pt,
  },
  subsection = {
    format = \zihao{4}\heiti,
    name = {,.},
    number = \arabic{section}.\arabic{subsection},
    beforeskip = 6pt,
    afterskip = 3pt,
  },
  subsubsection = {
    format = \zihao{-4}\heiti,
    beforeskip = 3pt,
    afterskip = 0pt,
  },
}

% 图题宋体五号居中（置于图下）；表题宋体五号加粗居中（置于表上）
\captionsetup[figure]{font=small,labelfont=small,labelsep=colon,position=below}
\captionsetup[table]{font={small,bf},labelfont={small,bf},labelsep=colon,position=above}
% 摘要标题：黑体小四加粗居中（ctexart 默认已居中，\abstractnamefont 并不存在，
% 所以直接通过 ctexset 把字体字号写进 abstractname）
\ctexset{abstractname = {\heiti\zihao{-4}摘要}}

% 正文：宋体小四，1.5 倍行距，首行缩进 2 字符
\linespread{1.5}
\setlength{\parindent}{2em}
\setlength{\parskip}{0pt}

\title{%%WP_TITLE%%}
\author{%%WP_AUTHOR%%}
\date{}

% 从用户模板抽取到的排版参数（覆盖上面的默认值；无则留空）
%%WP_STYLE_OVERRIDES%%

\begin{document}

\maketitle

%%WP_ABSTRACT_BEGIN%%
\begin{abstract}
%%WP_ABSTRACT%%
\keywords{%%WP_KEYWORDS%%}
\end{abstract}
%%WP_ABSTRACT_END%%

\newpage

%%WP_CONTENT%%

%%WP_BIBLIOGRAPHY%%

\end{document}
